\chapter{Mass bounds for slices of normal currents}\label{sec:normal-current-slicing}
Normal currents need not have rectifiable supports, so one cannot define a tangential derivative on such a support by invoking rectifiable charts. We instead prove the slicing mass bound first for smooth functions and then pass to Lipschitz functions through non-atomic levels. The argument develops \cite{Sim18}, Chapter 6, equations (4.6)--(4.11), pp.~184--185. All currents in this section have compact support in $\mathbb R^N$, and all flat norms have real coefficients.

\begin{lemma}[Integrated smooth slicing for a normal current]\label{lem:normal-slice-smooth-integrated}
Let $T$ be a normal $m$-current, $m\geq1$, and $f\in C^\infty(\mathbb R^N)$. Write $B_t=B_t^-$ for its lower distributional slice. For every test $(m-1)$-form $\omega$ and every $\psi\in C_c^\infty(\mathbb R)$,
\[
 \int\psi(t)B_t(\omega)\,dt=T(\psi(f)\,df\wedge\omega).
\]
Consequently the signed measure
\[
 A\longmapsto T(\chi_{f^{-1}(A)}df\wedge\omega)
\]
has Lebesgue density $B_t(\omega)$. This is the integrated form of \cite{Sim18}, equation (4.11).
\end{lemma}
\begin{proof}
Put $Q(s)=\int_s^\infty\psi(t)\,dt$. Fubini gives the difference $(T\llcorner Q(f))(d\omega)-((\partial T)\llcorner Q(f))(\omega)$. The smooth scalar product rule, Proposition \ref{prop:current-cutoff-boundary}, identifies it with $-T(d(Q(f))\wedge\omega)=T(\psi(f)df\wedge\omega)$.

The scalar function $B_t(\omega)$ is measurable and bounded by a fixed multiple of the two finite masses. It vanishes outside a bounded interval containing $f(\operatorname{spt}T)$: below that interval the restrictions vanish, while above it the two terms cancel. Thus it defines a finite signed measure. The right side is also a finite signed measure of the level variable. Equality on all smooth scalar tests implies equality as measures, by uniform density in continuous compactly supported tests and uniqueness in Riesz representation. This proves the final assertion without assuming a mass bound for $B_t$ in advance.
\end{proof}

\begin{proposition}[Smooth slicing bounds the integral of mass]\label{prop:normal-slice-smooth-mass}
Under the preceding assumptions, for every open $W\subset\mathbb R^N$ and Borel $A\subset\mathbb R$,
\[
 \int_A\mathbf M_W(B_t)\,dt
 \leq\int_{W\cap f^{-1}(A)}|Df|\,d\|T\|.
\]
The integrand is measurable, with infinity allowed before the estimate is proved. In particular $B_t$ has finite mass for almost every level. Compare \cite{Sim18}, equations (4.9)--(4.10).
\end{proposition}
\begin{proof}
Each $B_t$ has order at most one on every compact support, directly from its definition in terms of $T(d\omega)$ and $\partial T(\omega)$. Proposition \ref{prop:order-one-countable-mass} therefore provides a countable family $\{\omega_j\}_{j\geq1}$ of smooth $(m-1)$-forms compactly supported in $W$, each of comass at most one, with $\mathbf M_W(B_t)=\sup_j B_t(\omega_j)$. Include zero as the first test and include negatives. This test-space construction is independent of rectifiability.

For the first $k$ tests, partition $A$ into disjoint measurable sets $A_1,\ldots,A_k$ on which the corresponding test realizes their maximum, breaking ties by the smallest index. By Lemma \ref{lem:normal-slice-smooth-integrated},
\[
 \int_A\max_{1\leq j\leq k}B_t(\omega_j)\,dt
 =\sum_{j=1}^k T(\chi_{f^{-1}(A_j)}df\wedge\omega_j)
 \leq\int_{W\cap f^{-1}(A)}|Df|\,d\|T\|.
\]
Indeed the inverse images of the $A_j$ are disjoint, and $df\wedge\omega_j$ has comass at most $|Df|$. This bound holds for arbitrary polar $m$-vectors by mass/comass duality; simplicity is unnecessary. Ensure zero is the first test, so the maxima are nonnegative and increase. Monotone convergence as $k\to\infty$ gives the result and measurability. Taking $W=\mathbb R^N$ makes the right side finite because $T$ has compact support and $Df$ is bounded there.
\end{proof}

\begin{lemma}[Continuity under uniform perturbation of the slicing function]\label{lem:slice-function-uniform-limit}
Let $T$ be compactly supported and normal, and suppose $f_j,f$ are continuous near its support and $f_j\to f$ uniformly there. If
\[
 (\|T\|+\|\partial T\|)(\{f=t\})=0,
\]
then their lower distributional slices satisfy
\[
 \mathcal F_{\mathbb R}(B_t^{f_j}-B_t^f)\longrightarrow0.
\]
The lower slices are defined by their restrictions even when they have infinite mass. This justifies the approximation in the final paragraph of \cite{Sim18}, Chapter 6, Section 4.
\end{lemma}
\begin{proof}
Let $\delta_j=\sup_{\operatorname{spt}T}|f_j-f|$. The symmetric difference of the two strict sublevels is contained in $\{|f-t|\leq\delta_j\}$ on the current support. Subtraction of the defining formulas gives a flat decomposition with filling
$T\llcorner(\chi_{\{f_j<t\}}-\chi_{\{f<t\}})$
and residual the negative analogous restriction of $\partial T$. Their mass sum is at most
\[
 (\|T\|+\|\partial T\|)(\{|f-t|\leq\delta_j\}).
\]
This tends to zero by finiteness of the measure and the assumed zero mass on the level. The conclusion also holds when some $\delta_j=0$, by equality of the restrictions. Uniform approximation is used with respect to the current's support, rather than almost everywhere in ambient volume.
\end{proof}

\begin{theorem}[Lipschitz slicing of compact normal currents]\label{thm:normal-lipschitz-slicing}
Let $T$ be a compactly supported normal $m$-current in $\mathbb R^N$, $m\geq1$, and let $f:\mathbb R^N\to\mathbb R$ be $L$-Lipschitz. For almost every $t$, the common lower and upper slice $B_t$ is a compactly supported normal $(m-1)$-current. For every open $W$,
\[
 \int_{\mathbb R}\mathbf M_W(B_t)\,dt\leq L\|T\|(W).
\]
If $m\geq2$, then $\partial B_t=-B_t^{\partial T}$ at almost every level and
\[
 \int_{\mathbb R}\mathbf M_W(\partial B_t)\,dt
 \leq L\|\partial T\|(W).
\]
For $m=1$, the slice boundary is zero. No rectifiability or integer-multiplicity conclusion is asserted. These are the normal-current slicing estimates of \cite{Sim18}, equations (4.8)--(4.10), with compact support used to state global mass integrals.
\end{theorem}
\begin{proof}
Let $f_j$ be standard smooth ambient mollifications with radii tending to zero. Proposition \ref{prop:lipschitz-mollification} gives uniform convergence to $f$ and $|Df_j|\leq L$. The finite measure $f_\#(\|T\|+\|\partial T\|)$ has at most countably many atoms. At every other level, Lemma \ref{lem:slice-function-uniform-limit} gives flat, hence weak, convergence of the corresponding slices. For fixed open $W$, mass is lower semicontinuous under weak convergence: it is the supremum of the continuous evaluations on tests of comass at most one supported in $W$. Fatou and Proposition \ref{prop:normal-slice-smooth-mass} yield
\[
 \int\mathbf M_W(B_t^f)\,dt
 \leq\liminf_j\int\mathbf M_W(B_t^{f_j})\,dt
 \leq L\|T\|(W).
\]
Measurability follows from the same countable test family as in that proposition, applied to the order-one restriction formula. Taking $W=\mathbb R^N$ proves finite mass at almost every level. The support is compact by Proposition \ref{prop:all-level-slice-support-flat}.

When $m\geq2$, repeat the mass argument for the normal current $\partial T$, whose boundary is zero. At every level the distributional algebra gives $\partial B_t^T=-B_t^{\partial T}$. On the intersection of the two full-measure finite-mass sets this proves normality and the boundary integral estimate. In dimension one a zero-current has zero boundary, so the first mass conclusion already proves normality. Finally the two one-sided slices agree off the at most countable atomic set by Theorem \ref{thm:slice-flat-one-sided}. The quantitative measure inequality holds for each $W$; the full-measure set on which slices are normal is chosen once using their global masses.
\end{proof}

\begin{corollary}[Almost every restriction of a normal current is normal]\label{cor:normal-current-good-restriction}
Under Theorem \ref{thm:normal-lipschitz-slicing}, $T\llcorner\{f<t\}$ is normal for almost every $t$, with
\[
 \partial(T\llcorner\{f<t\})=(\partial T)\llcorner\{f<t\}+B_t.
\]
For every finite interval $(a,b)$ and every $\varepsilon>0$, a good level can be chosen in that interval, outside any prescribed null set, with
\[
 \mathbf M(B_t)\leq\frac{L\mathbf M(T)}{b-a}+\varepsilon.
\]
This is the normal-current analogue of the radius selection used for integral currents; see \cite{Sim18}, equation (4.10).
\end{corollary}
\begin{proof}
The restriction always has mass at most $\mathbf M(T)$. At a level of finite slice mass, its defining boundary formula has two finite-mass terms and gives normality. The integral of slice mass over $(a,b)$ is no larger than its full integral $L\mathbf M(T)$. If every good level outside the prescribed null set violated the stated bound, integration over the interval would contradict this inequality. The assertion uses a convenient global bound; it makes no claim of optimality for the chosen interval.
\end{proof}

\begin{corollary}[The mass bound uses only the intervening slab]\label{cor:normal-slice-interval-mass}
Under Theorem \ref{thm:normal-lipschitz-slicing}, for every open $W$ and $a<b$,
\[
 \int_a^b\mathbf M_W(B_t)\,dt
 \leq L\|T\|(W\cap\{a<f<b\}).
\]
For $m\geq2$, the corresponding boundary integral is at most $L\|\partial T\|(W\cap\{a<f<b\})$. Thus the good-level bound in Corollary \ref{cor:normal-current-good-restriction} can be sharpened to
$L\|T\|(\{a<f<b\})/(b-a)+\varepsilon$.
This is the interval estimate in \cite{Sim18}, equation (4.10).
\end{corollary}
\begin{proof}
The set $W'=W\cap\{a<f<b\}$ is open. Since $B_t$ is supported on $f^{-1}(t)$, its mass in $W'$ equals its mass in $W$ for $a<t<b$, and vanishes otherwise. Apply Theorem \ref{thm:normal-lipschitz-slicing} with $W'$ and integrate over all levels. The same support argument applies to the boundary slice. Averaging over $(a,b)$ proves the strengthened level-selection estimate. In particular the masses on the two endpoint levels are not added to the open-slab right side.
\end{proof}

\begin{proposition}[Smooth localization commutes with slicing]\label{prop:slice-smooth-localization}
Let $T$ be a locally normal $m$-current on $U$, $m\geq1$, let $f$ be locally Lipschitz, and let $\chi\in C_c^\infty(U)$. At every level, as distributional currents,
\[
 B_t^{\chi T}=\chi B_t^T.
\]
The same holds for upper slices. Here multiplication of an arbitrary distributional current by $\chi$ means evaluation on $\chi$ times the test form. This localization expands the local formulation of \cite{Sim18}, equations (4.6)--(4.10).
\end{proposition}
\begin{proof}
The Leibniz identity $\partial(\chi R)=\chi\partial R-R\llcorner d\chi$ holds distributionally by evaluating on a test form and expanding $d(\chi\omega)$; no finite mass of $\partial R$ is required. Apply it to $R=T\llcorner\{f<t\}$. On subtracting the restriction of $\partial(\chi T)=\chi\partial T-T\llcorner d\chi$, the two contractions cancel, because smooth contraction and Borel restriction of the finite-mass current $T$ commute. What remains is exactly $\chi B_t^T$. Closed sublevels give the upper-slice formula.
\end{proof}

\begin{theorem}[Local slicing of locally normal currents]\label{thm:local-normal-current-slicing}
Let $T$ be a locally normal $m$-current on $U\subset\mathbb R^N$, $m\geq1$, and let $f:U\to\mathbb R$ be locally Lipschitz. Outside a single Lebesgue-null set of levels, its slice is locally normal, is supported in $\operatorname{spt}T\cap f^{-1}(t)$, and satisfies
\[
 \partial B_t^T=-B_t^{\partial T}\quad(m\geq2).
\]
If $V\Subset U$ is open and $f$ is $L$-Lipschitz on an open neighborhood of $\overline V$, then
\[
 \int_{\mathbb R}\mathbf M_V(B_t^T)\,dt\leq L\|T\|(V).
\]
The analogous bound for its boundary holds with $\|\partial T\|$ when $m\geq2$. This is the local normal-current setting of \cite{Sim18}, pp.~184--185.
\end{theorem}
\begin{proof}
First fix such a neighborhood and choose a smooth cutoff $0\leq\chi\leq1$ supported there, equal to one near $\overline V$. The current $\chi T$, extended by zero to $\mathbb R^N$, is compactly supported and normal by the smooth product rule. Proposition \ref{prop:mcshane-extension} extends the restriction of $f$ on that neighborhood to a globally $L$-Lipschitz function $\widetilde f$; explicitly one can take the infimum of $f(y)+L|x-y|$ over points $y$ in the neighborhood. It equals $f$ there by its Lipschitz inequality.

Apply Theorem \ref{thm:normal-lipschitz-slicing} to $(\chi T,\widetilde f)$. Proposition \ref{prop:slice-smooth-localization} identifies its slice with that of $T$ on $V$. On this open set $\|\chi T\|=\|T\|$ and $\|\partial(\chi T)\|=\|\partial T\|$, because $\chi=1$ on a neighborhood. This proves the stated local mass estimates and normality on $V$ at almost every level.

Cover $U$ by countably many relatively compact balls on neighborhoods of which $f$ is Lipschitz. Remove the null sets furnished for those balls and, when needed, for $\partial T$. Local normality then holds simultaneously throughout $U$. The distributional boundary identity follows by $\partial^2=0$ exactly as in the compact case. The support statement is local: on the open strict sublevel the two terms cancel, on the open strict superlevel they vanish, and tests off $\operatorname{spt}T$ vanish. This gives the conclusion even when neither $T$ nor its boundary has finite total mass in $U$.
\end{proof}
